\begin{pastexam}[label={px:ncku-113-prog-02}]{成大資訊聯招 113 學年度 第 2 題}
（單選）

\textbf{[1\%, 1\%, 3\%]} Considering the data pairs with keys in the given
order: 20, 5, 10, 18, 4, 22, 11, 32, 21, as inputs to create a tree.
\begin{enumerate}[label=\roman*)]
\item If you create a \underline{max heap} for them, what is the level of the node containing 5?

(A)~1 \qquad (B)~2 \qquad (C)~3 \qquad (D)~4 \qquad (E)~None of the above
\item If you create a \underline{min heap} for them, which of the following statements is true?
\begin{enumerate}[label=(\Alph*)]
\item The node containing 11 is in the left subtree of the root.
\item The node containing 18 is in the right subtree of the root.
\item The node containing 20 is the root.
\item The node containing 21 is at level 3.
\item None of the above.
\end{enumerate}
\item If you create a \underline{symmetric min-max heap} for them, which of the following statements is true?
\begin{enumerate}[label=(\Alph*)]
\item To delete the node containing 5, the node containing 21 is used to replace the deleted node.
\item There are two nodes at level 4.
\item The leftmost grandchild of the root is 5.
\item The keys of the nodes in the right subtree of the root are 32, 10, and 20.
\item None of the above.
\end{enumerate}
\end{enumerate}
\end{pastexam}

\begin{pastsoln}
題目說 ``in the given order''，三小題都解讀成依序一個一個插入。 層數照題目的慣例（也是資結書的
約定）從根算第 $1$ 層。 答案：i) (D)；ii) (E)；iii) (A)。

i) 依序插入最大堆積，每次之後的陣列：
\begin{center}
\small
\begin{tabular}{l*{9}{c}}
\toprule
插入 & 1 & 2 & 3 & 4 & 5 & 6 & 7 & 8 & 9 \\
\midrule
20 & 20 \\
5 & 20 & 5 \\
10 & 20 & 5 & 10 \\
18 & 20 & 18 & 10 & 5 \\
4 & 20 & 18 & 10 & 5 & 4 \\
22 & 22 & 18 & 20 & 5 & 4 & 10 \\
11 & 22 & 18 & 20 & 5 & 4 & 10 & 11 \\
32 & 32 & 22 & 20 & 18 & 4 & 10 & 11 & 5 \\
21 & 32 & 22 & 20 & 21 & 4 & 10 & 11 & 5 & 18 \\
\bottomrule
\end{tabular}
\end{center}
插入 $32$ 時它從位置 $8$ 一路換到根，把 $5$ 帶到位置 $8$。 位置 $8$ 的深度是
$\lfloor \lg 8 \rfloor = 3$，在第 $4$ 層，選 (D)。 （如果題目的層數從 $0$ 算，答案會變成 (C)；成大
的考卷沿用從 $1$ 算的定義。）

ii) 依序插入最小堆積：
\begin{center}
\small
\begin{tabular}{l*{9}{c}}
\toprule
插入 & 1 & 2 & 3 & 4 & 5 & 6 & 7 & 8 & 9 \\
\midrule
20 & 20 \\
5 & 5 & 20 \\
10 & 5 & 20 & 10 \\
18 & 5 & 18 & 10 & 20 \\
4 & 4 & 5 & 10 & 20 & 18 \\
22 & 4 & 5 & 10 & 20 & 18 & 22 \\
11 & 4 & 5 & 10 & 20 & 18 & 22 & 11 \\
32 & 4 & 5 & 10 & 20 & 18 & 22 & 11 & 32 \\
21 & 4 & 5 & 10 & 20 & 18 & 22 & 11 & 32 & 21 \\
\bottomrule
\end{tabular}
\end{center}
結果是 $\langle 4,\allowbreak 5,\allowbreak 10,\allowbreak 20,\allowbreak 18,\allowbreak 22,\allowbreak 11,\allowbreak 32,\allowbreak 21 \rangle$。 根的左子樹是位置 $2,\allowbreak 4,\allowbreak 5,\allowbreak 8,\allowbreak 9$，也就是 $5,\allowbreak 20,\allowbreak 18,\allowbreak 32,\allowbreak 21$；右子樹是位
置 $3, 6, 7$，也就是 $10, 22, 11$。 (A) $11$ 在右子樹，錯。 (B) $18$ 在左子樹，錯。 (C) 根是
$4$，錯。 (D) $21$ 在位置 $9$，第 $4$ 層，錯。 所以選 (E)。

iii) 照資結書第~8.5.3~節的 \proc{SMMH-Insert}，下表是每次之後的 $A[2..]$：
\begin{center}
\small
\begin{tabular}{l*{9}{c}}
\toprule
插入 & 2 & 3 & 4 & 5 & 6 & 7 & 8 & 9 & 10 \\
\midrule
20 & 20 \\
5 & 5 & 20 \\
10 & 5 & 20 & 10 \\
18 & 5 & 20 & 10 & 18 \\
4 & 4 & 20 & 10 & 18 & 5 \\
22 & 4 & 22 & 10 & 18 & 5 & 20 \\
11 & 4 & 22 & 10 & 18 & 5 & 20 & 11 \\
32 & 4 & 32 & 10 & 22 & 5 & 20 & 11 & 18 \\
21 & 4 & 32 & 10 & 22 & 5 & 20 & 11 & 18 & 21 \\
\bottomrule
\end{tabular}
\end{center}
插入 $5$：位置 $3$ 是右子節點，比左兄弟 $20$ 小，對調。 插入 $4$：位置 $6$ 的祖父是根，
$4 < A[2] = 5$，$5$ 下來。 插入 $22$：比 $A[3] = 20$ 大，$20$ 下來。 插入 $32$：位置 $9$ 的祖父是
$2$，$32 > A[5] = 18$，$18$ 下來；位置 $5$ 的祖父是根，$32 > A[3] = 22$，$22$ 下來，$32$ 到位置
$3$。 插入 $21$：位置 $10$ 的祖父是 $2$，$A[4] = 10 \le 21 \le A[5] = 22$，不動。 最後的樹如下
圖 (c)：空的根，第 $2$ 層 $4, 32$，第 $3$ 層 $10, 22, 5, 20$，第 $4$ 層 $11, 18, 21$。

(B) 第 $4$ 層有三個節點，錯。 (C) 根的孫節點由左到右是 $10, 22, 5, 20$，最左邊是 $10$，錯。
(D) 根的右子樹是位置 $3, 6, 7$，鍵是 $32, 5, 20$，不是 $32, 10, 20$，錯。 (A)：資結書只定義了
取出最小值和最大值，題目沒有說刪除任意節點怎麼做；照取出最小值的做法類推，刪掉位置 $6$
的 $5$ 之後，拿最後一個節點的 $21$ 來補這個洞（洞沒有子節點；$21$ 比右兄弟 $A[7] = 20$ 大，
兩者再對調，最後是 $A[6] = 20$、$A[7] = 21$）。 「用 $21$ 那個節點來補」這件事成立，所以 (A)
對，選 (A)。 這也是唯一可能的答案：(B)、(C)、(D) 都確定是錯的。

\begin{center}
\begin{tikzpicture}[hn/.style={circle, draw=black!65, minimum size=0.66cm, inner sep=0pt, font=\small},
  ha/.style={hn, draw=memStroke, thick, fill=memFill},
  hb/.style={hn, draw=coreStroke, thick, fill=coreFill},
  he/.style={hn, draw=black!40, densely dashed},
  he-edge/.style={draw=black!55}]
\node[hn] (ya1) at (0.000,0.000) {32};
\node[hn] (ya2) at (-1.400,-0.950) {22};
\draw[he-edge] (ya1) -- (ya2);
\node[hn] (ya3) at (1.400,-0.950) {20};
\draw[he-edge] (ya1) -- (ya3);
\node[hn] (ya4) at (-2.100,-1.900) {21};
\draw[he-edge] (ya2) -- (ya4);
\node[hn] (ya5) at (-0.700,-1.900) {4};
\draw[he-edge] (ya2) -- (ya5);
\node[hn] (ya6) at (0.700,-1.900) {10};
\draw[he-edge] (ya3) -- (ya6);
\node[hn] (ya7) at (2.100,-1.900) {11};
\draw[he-edge] (ya3) -- (ya7);
\node[ha] (ya8) at (-2.450,-2.850) {5};
\draw[he-edge] (ya4) -- (ya8);
\node[hn] (ya9) at (-1.750,-2.850) {18};
\draw[he-edge] (ya4) -- (ya9);
\node[hn] (yb1) at (6.600,0.000) {4};
\node[hn] (yb2) at (5.200,-0.950) {5};
\draw[he-edge] (yb1) -- (yb2);
\node[hn] (yb3) at (8.000,-0.950) {10};
\draw[he-edge] (yb1) -- (yb3);
\node[hn] (yb4) at (4.500,-1.900) {20};
\draw[he-edge] (yb2) -- (yb4);
\node[hn] (yb5) at (5.900,-1.900) {18};
\draw[he-edge] (yb2) -- (yb5);
\node[hn] (yb6) at (7.300,-1.900) {22};
\draw[he-edge] (yb3) -- (yb6);
\node[hn] (yb7) at (8.700,-1.900) {11};
\draw[he-edge] (yb3) -- (yb7);
\node[hn] (yb8) at (4.150,-2.850) {32};
\draw[he-edge] (yb4) -- (yb8);
\node[hn] (yb9) at (4.850,-2.850) {21};
\draw[he-edge] (yb4) -- (yb9);
\node[font=\small] at (0,-3.65) {(a) 最大堆積，$5$ 在第 $4$ 層};
\node[font=\small] at (6.6,-3.65) {(b) 最小堆積};
\end{tikzpicture}

\begin{tikzpicture}[hn/.style={circle, draw=black!65, minimum size=0.66cm, inner sep=0pt, font=\small},
  ha/.style={hn, draw=memStroke, thick, fill=memFill},
  hb/.style={hn, draw=coreStroke, thick, fill=coreFill},
  he/.style={hn, draw=black!40, densely dashed},
  he-edge/.style={draw=black!55}]
\node[he] (yc1) at (0.000,0.000) {};
\node[hn] (yc2) at (-1.500,-0.950) {4};
\draw[he-edge] (yc1) -- (yc2);
\node[hn] (yc3) at (1.500,-0.950) {32};
\draw[he-edge] (yc1) -- (yc3);
\node[hn] (yc4) at (-2.250,-1.900) {10};
\draw[he-edge] (yc2) -- (yc4);
\node[hn] (yc5) at (-0.750,-1.900) {22};
\draw[he-edge] (yc2) -- (yc5);
\node[ha] (yc6) at (0.750,-1.900) {5};
\draw[he-edge] (yc3) -- (yc6);
\node[hn] (yc7) at (2.250,-1.900) {20};
\draw[he-edge] (yc3) -- (yc7);
\node[hn] (yc8) at (-2.625,-2.850) {11};
\draw[he-edge] (yc4) -- (yc8);
\node[hn] (yc9) at (-1.875,-2.850) {18};
\draw[he-edge] (yc4) -- (yc9);
\node[hn] (yc10) at (-1.125,-2.850) {21};
\draw[he-edge] (yc5) -- (yc10);
\node[font=\small] at (0,-3.65) {(c) 對稱最小最大堆積};
\end{tikzpicture}
\end{center}
\end{pastsoln}
